\documentclass[preview,border=2pt]{standalone}
\usepackage{listings}
\usepackage{xcolor}

\definecolor{darkgreen}{rgb}{0.0, 0.5, 0.0}
\lstdefinestyle{pythonstyle}{
    language=Python,
    basicstyle=\ttfamily\small,
    keywordstyle=\color{blue}\bfseries,
    stringstyle=\color{red},
    commentstyle=\color{darkgreen}\itshape,
    frame=single,
    framerule=0.5pt,
    breaklines=true,
    breakatwhitespace=false,
    showstringspaces=false,
    showtabs=false,
    tabsize=4,
    keepspaces=true,
    columns=fullflexible,
    upquote=true,
    extendedchars=true,
    inputencoding=utf8
}

\begin{document}
\begin{lstlisting}[style=pythonstyle]
import time
import rclpy
from rclpy.action import ActionServer
from rclpy.node import Node
from action_tutorials_interfaces.action import Fibonacci

class FibonacciActionServer(Node):
    def __init__(self):
        super().__init__('fibonacci_action_server')
        self._action_server = ActionServer(
            self, Fibonacci, 'fibonacci', self.execute_callback)

    def execute_callback(self, goal_handle):
        self.get_logger().info('Executando goal...')
        feedback_msg = Fibonacci.Feedback()
        feedback_msg.partial_sequence = [0, 1]

        for i in range(1, goal_handle.request.order):
            feedback_msg.partial_sequence.append(
                feedback_msg.partial_sequence[i] + feedback_msg.partial_sequence[i-1])
            goal_handle.publish_feedback(feedback_msg)
            time.sleep(1) # Simula tarefa longa

        goal_handle.succeed()
        result = Fibonacci.Result()
        result.sequence = feedback_msg.partial_sequence
        return result
\end{lstlisting}
\end{document}

